%% ---------------------------------------------------------------------------
%% proposal.tex
%%
%% Research Proposal, main document.
%%
%% ---------------------------------------------------------------------------
\documentclass[12pt,letterpaper]{article}
%\usepackage[english]{babel}     % supports english, but default is
\usepackage[spanish]{babel}
% include this if you want to import graphics files with /includegraphics
\usepackage{anysize}
\marginsize{2.5cm}{2.5cm}{1cm}{1cm}
\usepackage{textcomp}
\usepackage{url}
\bibliographystyle{unsrt}
\usepackage{graphics}
\usepackage{amssymb}
\usepackage{graphicx}
%\usepackage{slashbox}
\usepackage[latin1]{inputenc}
\usepackage{tikz}
% Color and strikethrough
\usepackage{color}
\usepackage{soul}

\usepackage{sectsty}
\allsectionsfont{\sffamily}

\definecolor{dblue}{RGB}{0,102,153}
\newcommand{\dB}[1]{\textcolor{dblue}{\textbf{#1}}}

\usepackage{longtable}
\usepackage{ifpdf}

\setlength{\parskip}{1em}

% Nombre del Estudiante
\newcommand{\scriptAuthor}{Daniel Moya S�nchez}

% T��tulo de la tesis
\newcommand{\scriptTitle}{Plan de Proyecto} 


% Keywords
\newcommand{\scriptKeywords}{key, words, ...}

% Para el PDF (cambiar si se desea otras cosas a lo indicado arriba
\newcommand{\pdfAuthor}{\scriptAuthor}
\newcommand{\pdfTitle}{\scriptTitle} 
\newcommand{\pdfKeywords}{\scriptKeywords}


\begin{document}
 
 \graphicspath{{./}{./fig/}}

 %% ---------------------------------------------------------------------------
%% titlepage.tex
%%
%% Title page
%%
%% ---------------------------------------------------------------------------

\thispagestyle{empty} 

\begin{center}

\textsc{\LARGE Instituto Tecnol\'ogico de Costa Rica} \\
\textsc{\Large \'Area Acad\'emica de Ingenier\'ia en Computadores}

\textsc{\Large Proyecto de Dise\~no en Ingenier\'ia en Computadores}


\par\vspace{20mm}

\includegraphics[scale=0.24]{logoTEC}

\par\vspace*{\fill}

{\LARGE\bf{\textsf{ \Huge \scriptTitle}}}

\par\vspace*{\fill}

Chair for Embedded System (CES) \\
Kalrsruhe Institute of Technology (KIT) \\
Period: 19/02/2018 (week 2) - 23/02/2018 (week 4) 
%Master Thesis {\sf Proposal} \\ 
%in fulfillment of the requirements for the degree of
%Plan de Proyecto
%Master of Science in Electronics Engineering \\
%Emphasis on Embedded Systems

\par\vspace{20mm}

\textsc{\Large \scriptAuthor}

\vspace*{\fill}

{\today}

\end{center}
\newpage 
\cleardoublepage  


%  \tableofcontents

 \clearpage

% El nombre del proyecto es preciso, conciso y no ambiguo
\section{Nombre del proyecto}

% Se identifica claramente la instituci�n donde se realizar� el proyecto
\section{Nombre de la instituci�n}

% Se identifica claramente qu� informaci�n se puede publicar, qu�
% informaci�n se puede compartir de manera limitada y qu�
% informaci�n se puede compartir si se firma un acuerdo de
% confidencialidad.
\section{Requerimientos de confidencialidad}


% El problema se describe de manera detallada, identificando los
% antecedentes, el contexto, usuarios, restricciones as� como otros
% aspectos relevantes para la resoluci�n del mismo.

% OPCIONAL: Se presenta una descripci�n del estado del arte de
% las investigaciones sobre el tema en cuesti�n.
\section{Descripci�n del problema}


\section{Objetivos}

% El objetivo general es directamente trazable al problema planteado, es claro, conciso.
\subsection{Objetivo general}

% Los objetivos espec�ficos son suficientes y consistentes con el
% objetivo general
\subsection{Objetivos espec�ficos}


% La identificaci�n de los involucrados y sus relaciones con el
% proyecto es detallada. Las relaciones se describen en detalle,
% se�alando cu�l es el inter�s del involucrado.
\section{Personas involucradas en el proyecto}


% Se presenta una descripci�n de la soluci�n propuesta en
% correspondencia con el problema planteado y los objetivos.
% Se utilizan diagramas de bloques u otra informaci�n que permita
% visualizar la estrategia de soluci�n al problema planteado.

% OPCIONAL: Se pressentan varias alternativas de soluci�n a estudiar.
\section{Descripci�n de la soluci�n}


% Los entregables con consistentes con los objetivos del proyecto y
% con los productos de un proyecto de ingenier�a de software para
% efectos de desarrollo posterior o mantenimiento.
% Los entregables tienen un identificador �nico.
% Los entregables deben incluir:
%	- Documento de Requerimientos
%	- Documento de Dise�o
%	- Plan de pruebas/Resultados de las preubas
%	- Otra documentaci�n relevante: instrucciones de compilaci�n,
% manuales de instalaci�n, manual de usuario etc.
\section{Entregables y criterios de aceptaci�n}


% Las cuatro �reas analizadas (personal, herramientas, procesos,
% insumos) son analizadas en detalle, identificando los riesgos m�s
% sobresalientes para cada una de esas �reas. Para cada riesgo se
% indica la probabilidad y el impacto para el proyecto.
\section{An�lisis de riesgos}

% Las actividades propuestas son consistentes con los objetivos
% espec�ficos.
% Las estimaciones de esfuerzos son proporcionales con la
% complejidad o el tama�o de las mismas.
% Se asigna un identificador �nico (c�digo) a cada actividad.
\section{Actividades y presupuesto de esfuerzo}

% El cronograma incluye todas las actividades y las mismas no
% guardan una consistecia l�gica para la construcci�n de los
% entregables.
% Se incluyen hitos para los entregables m�s sobresalientes.
% El nivel de granularidad de las actividades ofrece suficiente detalle.
% Al menos una actividad por semana.
\section{Cronograma}

% Referencias del Background y el Related Work
\bibliographystyle{sty/plainurl}
\bibliography{references}



\end{document}

